\documentclass[12pt, twoside]{article}
\input{../preamble}
\usepackage{pgfplots}
\pgfplotsset{compat=1.18}
\usepackage{eurosym}

\fancyhead[LE]{\thepage}
\fancyhead[RO]{\thepage \\ First \& last name: \hspace{2.25cm} \,\\ 11th Grade \hspace{2.9cm} \,}
\fancyhead[LO]{La Scuola d'Italia / Huson / IB Math: Linear Regression \\* 2 March 2026}

\begin{document}

\subsubsection*{2.7 Classwork: Linear Regression}
\begin{enumerate}[itemsep=0.1cm]

%--- Problem 1: Scatter plot, line of best fit, regression line, r ---
\item The table shows six paired observations of variables $x$ and $y$. \hfill [5 marks]

\begin{center}
\begin{tabular}{|c|cccccc|}
  \hline
  $x$ & 2 & 6 & 9 & 13 & 17 & 21 \\
  \hline
  $y$ & 5 & 15 & 11 & 14 & 19 & 22 \\
  \hline
\end{tabular}
\end{center}

A scatter diagram of the data is shown below.

\begin{center}
\begin{tikzpicture}
\begin{axis}[
  width=7.5cm, height=7.5cm,
  axis lines=left,
  xlabel={$x$},
  ylabel={$y$},
  xmin=0, xmax=25,
  ymin=0, ymax=25,
  xtick={0,5,...,25},
  ytick={0,5,...,25},
  grid=both,
  grid style={dashed, gray!30},
  tick label style={font=\small},
]
\addplot[only marks, mark=*, mark size=2pt] coordinates {
  (2,5) (6,15) (9,11) (13,14) (17,19) (21,22)
};
\end{axis}
\end{tikzpicture}
\end{center}

\begin{enumerate}[itemsep=0.25cm]
  \item On the scatter diagram above, draw a line of best fit. \hfill [1]

  \item Using your line, write down the approximate equation in the form $y = ax + b$. \hfill [1]

  \vspace{1.2cm}

  \item Use your GDC to find: \hfill [2]
  \begin{enumerate}[label=(\roman*), itemsep=0.5cm]
    \item the equation of the regression line of $y$ on $x$, giving $a$ and $b$ correct to
    three significant figures;
    \item the value of Pearson's product-moment correlation coefficient, $r$.
  \end{enumerate}

  \vspace{1.5cm}

  \item Use the guidance below to describe the correlation between $x$ and $y$. \hfill [1]

  \begin{quote}
    $0 \le |r| < 0.4$: weak, \quad
    $0.4 \le |r| < 0.8$: moderate, \quad
    $0.8 \le |r| \le 1$: strong.
  \end{quote}

  \vspace{0.8cm}

\end{enumerate}

\newpage

%--- Problem 2: Italian cities, latitude vs. salary ---
\item Eight Italian cities were surveyed to investigate the relationship between
latitude and average annual salary. The data are shown in the table below.
\hfill [7 marks]

\begin{center}
\begin{tabular}{|l|c|r|}
  \hline
  City & Latitude, $L$ ($^\circ$N) & Average salary, $S$ (\euro/year) \\
  \hline
  Milano  & 45.5 & 37,189 \\
  Torino  & 45.1 & 26,840 \\
  Bologna & 44.5 & 29,300 \\
  Genova  & 44.4 & 27,000 \\
  Firenze & 43.8 & 28,646 \\
  Roma    & 41.9 & 28,000 \\
  Napoli  & 40.9 & 24,000 \\
  Palermo & 38.1 & 22,000 \\
  \hline
\end{tabular}
\end{center}

\begin{enumerate}[itemsep=0.25cm]

  \item Find the value of Pearson's product-moment correlation coefficient, $r$.
  \hfill [1]

  \vspace{1.2cm}

  \item Use the guidance below to describe the correlation between latitude and
  average annual salary. \hfill [1]

  \begin{quote}
    $0 \le |r| < 0.4$: weak, \quad
    $0.4 \le |r| < 0.8$: moderate, \quad
    $0.8 \le |r| \le 1$: strong.
  \end{quote}

  \vspace{0.6cm}

  \item Write down the equation of the regression line of $S$ on $L$, in the
  form $S = aL + b$, where $a$ and $b$ are given to three significant figures.
  \hfill [2]

  \vspace{2cm}

  \item Bari is an Italian city at latitude $41.1^\circ$N. Use your regression
  line to estimate the average annual salary in Bari. \hfill [2]

  \vspace{2cm}

  \item Interpret the gradient of your regression line in the context of
  this problem. \hfill [1]

  \vspace{2cm}

\end{enumerate}

\end{enumerate}
\end{document}

Solutions
  Problem 1 — Linear regression, data: (2,5), (6,15), (9,11), (13,14), (17,19), (21,22)
  - (a) Line of best fit: should pass roughly through (0, 6) and (21, 22); accept any reasonable line
  - (b) Approximate equation: answers vary; e.g. y ≈ 0.8x + 5 (accept a wide range)
  - (c)(i) Regression line: y = 0.767x + 5.64  (a = 0.767, b = 5.64, both to 3 s.f.)
  - (c)(ii) r = 0.905 (3 s.f.)
  - (d) Strong positive correlation (0.8 ≤ r ≤ 1)

  Problem 2 — Italian cities: latitude (L, °N) vs. average annual salary (S, €/year)
  - (a) r = 0.746 (3 s.f.)
  - (b) Moderate positive correlation (0.4 ≤ 0.746 < 0.8)
       Note: Milan is an outlier — similar latitude to Turin but ~€10,000 higher salary,
       which weakens r relative to expectation.
  - (c) S = 1320L − 28800  (a = 1320, b = −28800, both to 3 s.f.)
  - (d) Using GDC stored values: S(41.1) ≈ €25,300
       Using 3 s.f. coefficients: S = 1320 × 41.1 − 28800 = €25,500 (accept either)
  - (e) For each additional degree of latitude (1° northward), average annual salary
        is predicted to increase by approximately €1320.
